\corrigesChapitre{10}{Puissances et suites matricielles}
\begin{corrige}{1}{Puissances diagonales}
\begin{questions}
\item $\begin{pmatrix}4&0\\0&9\end{pmatrix}$.
\item $\begin{pmatrix}2^n&0\\0&3^n\end{pmatrix}$ pour $n\in\N$.
\item $\begin{pmatrix}1/2&0\\0&1/3\end{pmatrix}$.
\end{questions}
\end{corrige}
\begin{corrige}{2}{Matrice périodique}
\begin{questions}
\item $I_2$.
\item $I_2$, car $2026$ est pair.
\item $A$, car $2027$ est impair.
\end{questions}
\end{corrige}
\begin{corrige}{3}{Premières colonnes}
\begin{questions}
\item $\binom11$.
\item $\binom21$.
\item $U_n=\binom n1$, confirmé par la formule de $A^n$.
\end{questions}
\end{corrige}
\begin{corrige}{4}{Équilibre}
\begin{questions}
\item $L=\tfrac12L+\binom12$, donc $L=\binom24$.
\item $U_n-L=2^{-n}(U_0-L)$.
\item La limite est $L$, car $2^{-n}\to0$.
\end{questions}
\end{corrige}
\begin{corrige}{5}{Prouver une puissance}
\begin{questions}
\item La formule donne $I_2$.
\item $\begin{pmatrix}1&n\\0&1\end{pmatrix}A=\begin{pmatrix}1&n+1\\0&1\end{pmatrix}$.
\item $\begin{pmatrix}1&100\\0&1\end{pmatrix}$.
\end{questions}
\end{corrige}
\begin{corrige}{6}{Une puissance triangulaire}
\begin{questions}
\item $\begin{pmatrix}9&5\\0&4\end{pmatrix}$.
\item $\begin{pmatrix}3^n&3^n-2^n\\0&2^n\end{pmatrix}$ ; pour $n=1$, le coefficient supérieur droit vaut $1$.
\item $\begin{pmatrix}243&211\\0&32\end{pmatrix}$.
\end{questions}
\end{corrige}
\begin{corrige}{7}{Ordre $3$}
\begin{questions}
\item $\begin{pmatrix}1&2&1\\0&1&2\\0&0&1\end{pmatrix}$.
\item $\begin{pmatrix}1&4&6\\0&1&4\\0&0&1\end{pmatrix}$.
\item $\begin{pmatrix}6\\4\\1\end{pmatrix}$, la troisième colonne de $T^4$.
\end{questions}
\end{corrige}
\begin{corrige}{8}{Deux récurrences}
\begin{questions}
\item $U_n=\binom{u_n}{v_n}$, $A=\operatorname{diag}(1/2,1/4)$, $C=\binom13$.
\item $L=\binom24$.
\item $u_n=2(1-2^{-n})\to2$ et $v_n=4(1-4^{-n})\to4$.
\end{questions}
\end{corrige}
\begin{corrige}{9}{Une méthode avec $PDP^{-1}$}
\begin{questions}
\item $P^{-1}=\begin{pmatrix}1&1\\0&1\end{pmatrix}$ ; le produit vaut $B=\begin{pmatrix}3&1\\0&2\end{pmatrix}$.
\item Les facteurs $P^{-1}P$ se simplifient dans les produits successifs ; une récurrence formalise l’identité.
\item $B^n=\begin{pmatrix}3^n&3^n-2^n\\0&2^n\end{pmatrix}$.
\end{questions}
\end{corrige}
\begin{corrige}{10}{Un couplage périodique}
\begin{questions}
\item $I_2$.
\item $U_{2k}=\binom10$, $U_{2k+1}=\binom0{1/2}$.
\item Non : les deux sous-suites constantes ont des valeurs différentes.
\end{questions}
\end{corrige}
\par\begingroup\renewcommand{\otBesoinPlace}[1]{}\noindent\begin{minipage}{\linewidth}
\begin{corrige}{11}{Point fixe non unique}
\begin{questions}
\item Toutes les colonnes $L$ satisfont $L=AL+C$.
\item $U_n=U_0$ pour tout $n$.
\item $I-A$ est la matrice nulle, non inversible. L’existence d’un point fixe ne garantit pas son unicité.
\end{questions}
\end{corrige}
\begin{corrige}{12}{Contrôler une formule en ordre $3$}
\begin{questions}
\item Pour $n=0$, les hors-diagonales sont nulles ; pour $n=1$, $(1,3)=0$ : on retrouve $I$ puis $T$.
\item $n+n(n-1)/2=n(n+1)/2$.
\item Le coefficient satisfait la récurrence $c_{n+1}=c_n+n$, avec $c_0=0$ ; c’est la somme $0+1+\cdots+(n-1)=\binom n2$.
\end{questions}
\end{corrige}
\end{minipage}\endgroup
